\documentclass[11pt,a4paper]{article}
\usepackage{fontspec,amsmath,amssymb,amsthm,microtype}
\setmainfont{lmroman10-regular.otf}[BoldFont=lmroman10-bold.otf,ItalicFont=lmroman10-italic.otf,BoldItalicFont=lmroman10-bolditalic.otf]
\setmonofont{lmmono10-regular.otf}[Scale=MatchLowercase]
\usepackage[margin=25mm]{geometry}
\usepackage[unicode,hidelinks]{hyperref}
\setlength{\emergencystretch}{3em}
\setlength{\parskip}{0.4em}
\setlength{\parindent}{0pt}
\allowdisplaybreaks[2]
\raggedbottom
\title{\texorpdfstring{Proposition 2 with actual Gaussian and Gibbs moments}{Proposition 2 with actual Gaussian and Gibbs moments}}
\author{LangevinLean mathematical exposition}
\date{}
\begin{document}\maketitle
\noindent Source module: \texttt{Kinetic\allowbreak{}Langevin/\allowbreak{}Second\allowbreak{}Order\allowbreak{}Lp\allowbreak{}Local.\allowbreak{}lean}. The accompanying HTML pairs this mathematical exposition with the exact Lean statements and records the scope of the formalization.


This theorem combines the actual random-Hessian Taylor estimate with the genuine cubic/\allowbreak{}quartic defect decomposition. Its reference is the actual coupled L\ensuremath{{}^{2}} solution: exact mild and kinematic equations, with each time marginal equal to the Gaussian-times-Gibbs target. The moment and local-error inequalities are derived inside the proof.


\section{\texorpdfstring{Derive all local constants from target marginals}{Derive all local constants from target marginals}}


Project the actual phase pushforward law to obtain the standard Gaussian velocity marginal and the Gibbs position marginal. The Gaussian fourth moment gives the displacement Taylor forcing bound M\ensuremath{{}_{2}}\ensuremath{\surd}(p\ensuremath{{}^{2}}+2p)t\ensuremath{{}^{2}}/\allowbreak{}2. The Gibbs score theorem gives \ensuremath{\Vert}\ensuremath{\nabla}f(X(t))\ensuremath{\Vert}\ensuremath{{}_{2}}\ensuremath{\le}\ensuremath{\surd}(Mp), while the actual frozen Hessian operator has norm at most M. The generic local decomposition therefore has constant at most M\ensuremath{{}_{2}}\ensuremath{\surd}(p\ensuremath{{}^{2}}+2p)+M\ensuremath{\surd}(Mp). It yields velocity error Ct\ensuremath{{}^{3}}/\allowbreak{}6 and position error Ct\ensuremath{{}^{4}}/\allowbreak{}24. If \ensuremath{\gamma}t\ensuremath{\le}1/\allowbreak{}5, the actual transformed phase error is at most Ct\ensuremath{{}^{3}}/\allowbreak{}4. The conclusions refer to the corrected paths retaining the entire free-OU response; separate kernel identities identify those paths with the implemented KLMC2 formula.



{\raggedright\noindent Lean declarations: \texttt{lp\_\allowbreak{}second\allowbreak{}Order\_\allowbreak{}stationary\_\allowbreak{}local\_\allowbreak{}bound}.\par}

\end{document}
